\section{The complete cubic directional layer of Tornheim zeta}
\label{sec:rays}

Define the uncolored Mordell--Tornheim function in its absolute
convergence region by
\begin{equation}\label{ray:T}
 \T(A,B,C)=\sum_{m,n\ge1}\frac1{m^A n^B(m+n)^C}.
\end{equation}
For instance, sufficiently large positive real parts suffice.
Its meromorphic continuation is symmetric in $A,B$, but is not
jointly holomorphic at the origin. We form a restriction away from
its polar hyperplanes before taking a one-variable limit.

Let
\[
 W_{a,b,c}(t)=\T(at,bt,ct),\qquad
 \Omega_3=\left.\frac{\mathrm d^3}{\mathrm dt^3}\T(t,t,t)\right|_{t=0}.
\]
The two admissibility conditions are
\begin{equation}\label{ray:admissible}
 (a+c)(b+c)\ne0.
\end{equation}
They allow complex and unequal slopes, including negative slopes.

\subsection{Statement of the formula}

Put
\begin{align}
 \sigma&=a+b+c,\notag\\
 \mathcal B&=\frac{c(c^2+ab-a^2-b^2)}{(a+c)(b+c)},\label{ray:B}\\
 \Delta&=a^2+b^2-c(a+b),\notag\\
 J_3&=W_{1,0,2}'''(0),\notag\\
 \kappa&=J_3+\frac{35}{2}\zeta'''(0)
                 +6L\zeta''(0)-9\zeta'''(-1).\label{ray:kappa}
\end{align}
The constant $J_3$ is a fixed Euler double-zeta directional derivative:
$\T(t,0,2t)=Z_2(2t,t;1)$. An independent convergent integral for
$\kappa$ is given below.

\begin{theorem}[All admissible cubic directions]\label{thm:cubic-rays}
Under \eqref{ray:admissible}, $W_{a,b,c}$ is holomorphic at zero and
\begin{align}
 W_{a,b,c}'''(0)={}&abc\,\Omega_3-\frac{c\Delta}{2}\kappa\notag\\
 &+\left(8abc-\frac{(a+c)^3+(b+c)^3}{2}\right)\zeta'''(0)\notag\\
 &-\frac32\bigl((a+b)(ab+c^2)-4abc\bigr)L\zeta''(0)
       +\sigma^2\mathcal B\,\zeta'''(-1).
 \label{ray:main}
\end{align}
Thus every third directional derivative belongs to the span of
$\Omega_3,\kappa,\zeta'''(0),L\zeta''(0),\zeta'''(-1)$ over
the field of rational functions of the slopes. By the cubic bridge,
$\Omega_3$ can be replaced by $\eta(1)$ and explicit ordinary zeta
and Stieltjes data.
\end{theorem}

Every coefficient is homogeneous of degree three. This agrees with
$W_{ra,rb,rc}(t)=W_{a,b,c}(rt)$. The formula contains no surviving
Euler-constant or Gamma-expansion coefficients in addition to those
already packaged into $\Omega_3$ and $\kappa$.

\subsection{A holomorphic remainder with exactly two visible polar divisors}

For $x>0$ set $F(A,B;x)=\Li_A(e^{-x})\Li_B(e^{-x})$.
Jonqui\`ere's formula gives, locally uniformly for $A$ near zero,
\begin{equation}\label{ray:Jonquiere}
 \Li_A(e^{-x})=\Gamma(1-A)x^{A-1}
       +\sum_{j\ge0}\zeta(A-j)\frac{(-x)^j}{j!},\qquad 0<x<2\pi.
\end{equation}
Subtract the six local terms
\begin{align}
 S(A,B;x)={}&\Gamma(1-A)\Gamma(1-B)x^{A+B-2}
 +\Gamma(1-A)\zeta(B)x^{A-1}\notag\\
 &+\Gamma(1-B)\zeta(A)x^{B-1}
 -\Gamma(1-A)\zeta(B-1)x^A\notag\\
 &-\Gamma(1-B)\zeta(A-1)x^B+\zeta(A)\zeta(B).
 \label{ray:S}
\end{align}
Define
\begin{align}
 H(A,B,C)={}&\int_0^1x^{C-1}\bigl(F(A,B;x)-S(A,B;x)\bigr)\dd x
       +\int_1^\infty x^{C-1}F(A,B;x)\dd x\notag\\
 &+\frac{\Gamma(1-A)\Gamma(1-B)}{A+B+C-2}
 +\frac{\Gamma(1-A)\zeta(B)}{A+C-1}
 +\frac{\Gamma(1-B)\zeta(A)}{B+C-1}.
 \label{ray:H}
\end{align}
All denominators on the last line are nonzero near the origin.
On a sufficiently small closed parameter polydisc, the first integrand
and every fixed parameter derivative are bounded at zero by
$x^{-\delta}(1+|\log x|^N)$ for a $\delta<1$; the second integral
has exponential decay. Thus $H$ is jointly holomorphic and symmetric
in $A,B$. Termwise integration and meromorphic continuation yield
\begin{equation}\label{ray:decomposition}
 \begin{split}
 \T(A,B,C)=\frac1{\Gamma(1+C)}\biggl\{&\zeta(A)\zeta(B)
 -\frac{C\Gamma(1-A)\zeta(B-1)}{A+C}\\
 &-\frac{C\Gamma(1-B)\zeta(A-1)}{B+C}
 +CH(A,B,C)\biggr\}.
 \end{split}
\end{equation}
This also proves removability on every admissible ray. The six-term
subtraction and the decomposition are retained from the incoming
coincident report \cite{IncomingCoincident}; the cubic determination
below is the additional step.

\subsection{Exact slices determine the constant and linear remainder}

Two elementary restrictions provide all the lower coefficients:
\begin{align}
 \T(0,0,t)&=\zeta(t-1)-\zeta(t),\label{ray:axis}\\
 \T(t,0,t)&=\frac{\zeta(t)^2-\zeta(2t)}2.\label{ray:symmetric-slice}
\end{align}
For the first, count positive pairs with $m+n=k$. For the second,
use $\T(A,0,C)=Z_2(C,A;1)$ and the exact stuffle product
\begin{equation}\label{ray:stuffle-slice}
 \T(A,0,C)+\T(C,0,A)=\zeta(A)\zeta(C)-\zeta(A+C).
\end{equation}
Both identities are first proved in domains of absolute convergence
and then extended meromorphically. A fixed $C=0$ substitution alone
would not determine the normal derivative of $CH$; this is why
\eqref{ray:symmetric-slice} supplies real additional information.

Write
\begin{align}
 H(A,B,C)={}&h_0+h_A(A+B)+h_C\,C\notag\\
 &+h_{20}(A^2+B^2)+h_{11}AB+h_{10c}(A+B)C+h_{02}C^2
 +O\bigl((|A|+|B|+|C|)^3\bigr).
 \label{ray:H-Taylor}
\end{align}
Here $h_C$ denotes the coefficient of the linear term in $C$;
$h_{02}$ is its quadratic coefficient. Direct comparison in
\eqref{ray:axis}--\eqref{ray:symmetric-slice} gives
\begin{align}
 h_0={}&\zeta'(-1)+\frac L2-\frac{5\gamma}{12},\label{ray:h0}\\
 h_C={}&\frac{\zeta''(-1)-\zeta''(0)}2
       -\gamma\left(\zeta'(-1)+\frac L2\right)
       +\frac5{24}(\gamma^2+\zeta(2)),\label{ray:hC}\\
 h_A={}&\frac{L^2}{8}-\frac{\zeta''(0)}2+
       \gamma\zeta'(-1)-\frac{\gamma L}{4}
       -\frac{\gamma^2+\zeta(2)}{24}.
 \label{ray:hA}
\end{align}
For example, the axis gives
$\Gamma(1+t)(\zeta(t-1)-\zeta(t))=5/12+tH(0,0,t)$,
which proves the first two lines. On the symmetric slice the
coefficient of $t^2$ then gives the third line.

Substitution now proves the complete lower layer
\begin{align}
 W(0)&=\frac14+\frac{c}{12(a+c)}+\frac{c}{12(b+c)},\notag\\
 W'(0)&=\frac{a+b+2c}{4}L+\mathcal B\zeta'(-1),\notag\\
 W''(0)&=\frac{2ab+c(a+b)}4L^2
 -\frac{a^2+b^2+2c(a+b)+2c^2}{2}\zeta''(0)
 +\sigma\mathcal B\zeta''(-1).
 \label{ray:lower}
\end{align}
These agree with the earlier report. In particular,
 $\omega(0)=1/3$, $\omega'(0)=L$, and
$\omega''(0)=L^2-4\zeta''(0)$ follow here without using the
unknown cubic moment.

\subsection{Proof of the cubic formula}

At quadratic degree, the symmetric slice fixes
$h_{20}+h_{10c}$, not the two terms separately. This distinction is
the main structural point. The axis fixes $h_{02}$; the diagonal
then fixes $h_{11}$, leaving exactly one more coefficient, determined
by $J_3$.

Here are explicit finite coefficient formulas which also make the
algebra completely reproducible. Set $V=h_{20}+h_{10c}$ and define
the cubic Taylor polynomials
\begin{align*}
 O(t)&=\frac13+Lt+\frac{L^2-4\zeta''(0)}2t^2
                         +\frac{\Omega_3}{6}t^3,\\
 D(t)&=\frac7{18}+\left(\frac54L+\zeta'(-1)\right)t
       +\left(\frac{L^2}{4}-\frac{13}{4}\zeta''(0)
                               +\frac32\zeta''(-1)\right)t^2
       +\frac{J_3}{6}t^3.
\end{align*}
Comparison on the four slices gives
\begin{align}
 h_{02}&=[t^3]\Gamma(1+t)(\zeta(t-1)-\zeta(t)),\notag\\
 V+h_{02}&=[t^3]\left\{
 \frac{\Gamma(1+t)}2(\zeta(t)^2-\zeta(2t))
 +\frac{\zeta(t)}2-\frac{\Gamma(1-t)}{24}+\zeta(t-1)\right\},\notag\\
 h_{11}+2V+h_{02}&=[t^3]\left\{
 \Gamma(1+t)O(t)-\zeta(t)^2+\Gamma(1-t)\zeta(t-1)\right\},\notag\\
 -2h_{20}+4V+8h_{02}&=[t^3]\left\{
 \Gamma(1+2t)D(t)+\frac{\zeta(t)}2
                      -\frac{\Gamma(1-t)}{18}+\zeta(t-1)\right\}.
 \label{ray:finite-system}
\end{align}
The coefficient system is triangular with nonzero diagonal. Thus
it determines all four quadratic coefficients for any assigned
$\Omega_3,J_3$.

Use
\[
 \Gamma(1\pm t)=1\mp\gamma t+
 \frac{\gamma^2+\zeta(2)}2t^2
 \mp\frac{\gamma^3+3\gamma\zeta(2)+2\zeta(3)}6t^3+O(t^4)
\]
and the ordinary Taylor expansions of $\zeta$ at $0$ and $-1$.
Substitute \eqref{ray:finite-system} into
\eqref{ray:decomposition}. The coefficients of $\Omega_3$ and
$J_3$ are respectively $abc$ and $-c\Delta/2$.
The remaining nonzero coefficients are
\begin{align*}
 [\zeta'''(0)]\colon\;&8abc-\frac{(a+c)^3+(b+c)^3}{2}
                                      -\frac{35}{4}c\Delta,\\
 [L\zeta''(0)]\colon\;&-\frac32\bigl((a+b)(ab+c^2)-4abc\bigr)
                                      -3c\Delta,\\
 [\zeta'''(-1)]\colon\;&\sigma^2\mathcal B+\frac92c\Delta.
\end{align*}
All coefficients involving $\gamma,\zeta(2),\zeta(3)$ or
$\zeta'(-1),\zeta''(-1)$ cancel. Absorbing the three multiples
of $c\Delta$ into \eqref{ray:kappa} proves
Theorem~\ref{thm:cubic-rays}. This is a finite algebraic proof;
the delivered symbolic certificate independently verifies the exact
rational identity with all constants treated as formal symbols.

\subsection{An independent convergent integral for the second coordinate}

Define
\begin{align}
 Q_2(x)&=\left.\partial_s^2\Li_s(e^{-x})\right|_{s=0}
               =\sum_{n\ge1}(\log n)^2e^{-nx},\notag\\
 P(x)&=(\log x+\gamma)^2+\zeta(2),\notag\\
 S_2(x)&=\left(x^{-2}-\frac1{2x}+\frac1{12}\right)P(x)
             +\zeta''(0)\left(x^{-1}-\frac12\right)-\zeta''(-1),\notag\\
 I_2&=\int_0^1\left(\frac{Q_2(x)}{e^x-1}-S_2(x)\right)\frac{\dd x}{x}
       +\int_1^\infty\frac{Q_2(x)}{x(e^x-1)}\dd x.
 \label{ray:I2}
\end{align}
Both integrals converge ordinarily: the first integrand is
$O(1+\log^2x)$; the second is $O(e^{-3x}/x)$.

\begin{proposition}[Mellin coordinate]\label{prop:kappa-integral}
The quadratic coefficient $h_{20}$ satisfies
\begin{equation}\label{ray:h20-integral}
 h_{20}=\frac\gamma4+\frac38-\frac{\zeta''(0)}2+\frac{I_2}{2},
\end{equation}
and
\begin{align}
 \kappa={}&-12h_{20}-\frac{\gamma^3}{6}
       -\frac{\gamma\zeta(2)}2+3\gamma\zeta''(0)
       +6\gamma\zeta''(-1)-\frac{\zeta(3)}3
       -3\zeta'''(0)-2\zeta'''(-1).
 \label{ray:kappa-h20}
\end{align}
\end{proposition}
\begin{proof}
Differentiate \eqref{ray:H} twice in $A$ at $(0,0,0)$.
The product derivative is $Q_2(x)/(e^x-1)$, and the derivative of
the six subtractions is $S_2(x)$. Differentiation under the subtracted
integrals follows from the local bounds used above. The rational
last line of \eqref{ray:H}, on $B=C=0$, is
\[
 \Gamma(1-A)\left(\frac1{A-2}-\frac1{2(A-1)}\right)-\zeta(A).
\]
Its coefficient of $A^2$ is
$\gamma/4+3/8-\zeta''(0)/2$, proving
\eqref{ray:h20-integral}. Expanding the last line of
\eqref{ray:finite-system} and using \eqref{ray:kappa} gives
\eqref{ray:kappa-h20}.
\end{proof}

The independent integral evaluation gives
\begin{align*}
 I_2&=-0.2860481161626797704387366252971602\ldots,\\
 \kappa&=-3.0226121641572464811701683702659737\ldots.
\end{align*}
These displayed digits are numerical diagnostics, not an arithmetic
characterization. The definition \eqref{ray:I2} supplies such a
characterization without referring back to an unknown generic ray.

\subsection{Coordinate cancellation and exact examples}

In the formal coordinate system of Theorem~\ref{thm:cubic-rays},
the $\kappa$ term disappears exactly when
\begin{equation}\label{ray:conic}
 c\bigl(a^2+b^2-c(a+b)\bigr)=0.
\end{equation}
For positive $a,b$, the choice
$c=(a^2+b^2)/(a+b)$ gives a continuum of asymmetric directions with
only the harmonic coordinate $\eta(1)$ remaining. For fixed real
$c$, the nontrivial locus can equivalently be written
$(a-c/2)^2+(b-c/2)^2=c^2/2$.

The $\Omega_3$ term disappears when $abc=0$. With $c\ne0$, both
displayed nonclassical coefficients vanish only on the axis
$(0,0,c)$ and the rays $(c,0,c)$ or $(0,c,c)$. With $c=0$,
admissibility requires $ab\ne0$ and the restriction equals
$\zeta(at)\zeta(bt)$. These are coefficient-elimination statements,
not assertions that other numerical cancellations are impossible.

Some explicit specializations are
\begin{align}
 W_{1,2,3}'''(0)={}&6\Omega_3+6\kappa
       -\frac{93}{2}\zeta'''(0)-\frac{27}{2}L\zeta''(0)
       +\frac{162}{5}\zeta'''(-1),\label{ray:123}\\
 W_{2,3,1}'''(0)={}&6\Omega_3-4\kappa
       +\frac52\zeta'''(0)-\frac{33}{2}L\zeta''(0)
       -18\zeta'''(-1),\label{ray:231}\\
 W_{1,1,2}'''(0)={}&2\Omega_3+2\kappa
       -11\zeta'''(0)-3L\zeta''(0)+\frac{32}{3}\zeta'''(-1),\notag\\
 W_{0,0,1}'''(0)={}&\zeta'''(-1)-\zeta'''(0),\notag\\
 W_{1,0,1}'''(0)={}&-\frac92\zeta'''(0)-\frac32L\zeta''(0).
\end{align}
The last two also follow directly from the exact slices.

\begin{remark}[Why two coordinates are sufficient, and what this does not prove]
Modulo the already determined constant and linear coefficients of
$H$, the axis and symmetric slice leave two formal quadratic degrees
of freedom. The diagonal and $(1,0,2)$ slices fix them separately.
Their contributions to generic cubic rays are independent polynomials
$abc$ and $c\Delta$. This proves uniqueness in the specified formal
reconstruction problem. It does not prove that $\eta(1)$ and $\kappa$
are independent over any arithmetic constant field.
\end{remark}
